
%%%%    TO INCLUDE/EXCLUDE THE SOLUTIONS, 
%       COMMENT OUT AT THE END OF EACH QUESTION.

\documentclass[a4paper,11pt]{article}
\usepackage[utf8]{inputenc}
\usepackage[english]{babel}
\usepackage[margin=3cm]{geometry}
%\usepackage[onehalfspacing]{setspace}
\usepackage{amsmath,graphicx}
\graphicspath{{figures/}}
\usepackage{microtype}
\usepackage{amssymb}
\usepackage{xcolor}
\usepackage{subcaption}
\usepackage[hidelinks]{hyperref}
%\usepackage{chngcntr}
\usepackage[justification=centering]{caption}
\setlength{\parindent}{0em}
\usepackage{hyperref}
\hypersetup{
    colorlinks=true,
    linkcolor=blue,
    filecolor=magenta,      
    urlcolor=blue,
    }
\begin{document}



\title{Advanced Macroeconomics (cod. 20285) \\ \textbf{Problem set 4}}
% \title{Advanced Macroeconomics (cod. 20285) \\ Solutions to Problem set 4}

\date{A.Y. 2025--26}

\author{Instructors: Antonella Trigari\\
Teaching Assistant: Henri Graul}
\maketitle
\begin{itemize}
\item Due date: May $3^{rd}$, by \textbf{23:59} CEST;

\item Write your solutions with a word processor (\LaTeX$\,$would be preferred);

\item \textbf{Grading} is based on \textbf{either Exercise 1 or 2} (randomly chosen) and \textbf{Exercise 3};

\item \textbf{Optional} questions will \textbf{not be graded}. They serve as additional practice exercises; 

\item Please submit your solutions to \href{mailto:henri.graul@unibocconi.it}{henri.graul@unibocconi.it};

\item Submit only \textbf{one} pdf for each group, don't forget to specify the names of the members of the group on the first page. Please also \textbf{indicate your names in the document title}, i.e. for example "PS4Graul/Trigari.pdf".     
\end{itemize}

\bigskip

	

%%%%%%%%%%%%%%%%%%%%%%%%%%%%%%%%%%%%%%%%%%%
%%%%%%%%%%%%%%%%%%%%%%%%%%%%%%%%%%%%%%%%%%%

% \include{2025/PS/PS 4/ex3_sol}

%%%%%%%%%%%%%%%%%%%%%%%%%%%%%%%%%%%%%%%%%%%
%%%%%%%%%%%%%%%%%%%%%%%%%%%%%%%%%%%%%%%%%%%
%%%%%%%%%%%%%%%%%%%%%%%%%%%%%%%%%%%%%%%%%%%

\newpage


\section*{OPTIONAL EXERCISES}

Note that the notation in these exercises might differ slightly from that used in the lectures.


%%%%%%%%%%%%%%%%%%%%%%%%%%%%%%%%%%%%%%%%%%%
%%%%%%%%%%%%%%%%%%%%%%%%%%%%%%%%%%%%%%%%%%%
%%%%%%%%%%%%%%%%%%%%%%%%%%%%%%%%%%%%%%%%%%%


\section*{Exercise 4: Optimal monetary policy under Price Indexation}

In this exercise, we will analyse how partial price indexation by firms impacts optimal monetary policy design. In this case, the second-order approximation to the household's welfare losses takes the form:
\begin{equation*}
    \frac{1}{2} \mathbb{E}_0 \left[ \sum_{t=0}^{\infty} \beta^t (\vartheta x_t^2 + \{ \pi_t  - \zeta \pi_{t-1} \}^2 )\right],
\end{equation*}
where $\zeta$ denotes the degree of price indexation to past inflation.

The equation describing the evolution of inflation is now given by:
\begin{equation*}
    \pi_t - \zeta \pi_{t-1} = \kappa x_t + \beta \mathbb{E}_t [\pi_{t+1} - \zeta \pi_t] + u_t,
\end{equation*}
where $u_t$ represents an exogenous i.i.d cost-push shock.

\begin{enumerate}
    \item Briefly explain the relationship between $x_t, \widetilde{y}_t$ and $u_t$. Why are we using $x_t$ when studying optimal monetary policy? Next, compare the Phillips curve in this exercise with the one you obtained in Exercise 1.5. Under what condition are they congruent?
    
    \item Determine the optimal policy under discretion. Specifically, we assume that for simplicity, the monetary authority seeks to minimize each
period \textit{only} the short-term losses, $\vartheta x_t^2 + \{ \pi_t  - \zeta \pi_{t-1} \}^2$. This implies that the central bank does not set-up a dynamic problem, given that past inflation is a state variable in this setting. Hence it faces the following constrained optimization problem: 
\begin{equation*}
    \min_{x_t, \pi_t} \vartheta x_t^2 + \{ \pi_t  - \zeta \pi_{t-1} \}^2
\end{equation*}
s.t.
$$\mathrm{s.t.}\quad\begin{cases}
\pi_t - \zeta \pi_{t-1} = \kappa x_t + \beta \mathbb{E}_t [\pi_{t+1} - \zeta \pi_t] + u_t \\
\pi_{t-1} \;\text{given}
\end{cases} $$



\item Determine the optimal policy under commitment. 
\item Discuss how the degree of indexation $\zeta$ affects the optimal responses
to a transitory cost-push shock under the previous two scenarios.
\item \textit{[Extra Optional:} Determine optimal policy under discretion \textit{without} the simplifying assumption made in 2.2.]
\end{enumerate}



%%%%%%%%%%%%%%%%%%%%%%%%%%%%%%%%%%%%%%%%%%%
%%%%%%%%%%%%%%%%%%%%%%%%%%%%%%%%%%%%%%%%%%%

%\include{2025/PS/PS 4/ex4_sol}

%%%%%%%%%%%%%%%%%%%%%%%%%%%%%%%%%%%%%%%%%%%
%%%%%%%%%%%%%%%%%%%%%%%%%%%%%%%%%%%%%%%%%%%



\section*{Exercise 5: Government Spending Shocks and Optimal Monetary Policy}

Consider an economy in which the government has a time-varying \textit{exogenous} share of purchases in output $\gamma _{t}$:

\begin{equation*}
G_{t}=\gamma _{t}Y_{t}, 
\end{equation*}
where $\gamma$ evolves exogenously, with $\gamma_t<1$.


The economy's resource constraint reads:
\begin{equation*}
Y_{t}=C_t+G_{t}.
\end{equation*}
The production function reads:
\begin{equation*}
Y_{t}=A_{t}N_{t},
\end{equation*}%
where $A_{t}$ is labor productivity, evolving exogenously as usual.

The representative household maximizes the intertemporal utility

\begin{equation*}
\mathbb{E}_{0} \sum_{t=0}^{\infty }\beta ^{t}\left \{ \frac{%
C_{t}^{1-\sigma }}{1-\sigma }-\frac{N_{t}^{1+\varphi }}{1+\varphi }
\right \}
\end{equation*}%
under a standard sequence of budget constraints. Assume $\sigma \geq 1$ and $\varphi \geq 0$.



\begin{enumerate}
\item[1.]  Formulate the social planner’s problem and derive the first-order conditions for efficiency. Derive an expression for the efficient level of output, denoted as $Y_{t}^{e}$. What are its determinants? In particular, what is the effect of a rise in $\gamma_{t}$? 

\item [2.] Assume now that the economy is characterized by imperfect competition and price stickiness à la Calvo, and define $\varepsilon$ as the elasticity of substitution between varieties. Write the expression for the marginal cost and use it, along with the price-setting condition from profit maximization and the labor supply condition from household optimization, to derive the expression for the natural level of output, $Y^n_t$. What are the determinants of $Y^n_t$? What is the effect of a rise in $\gamma_{t}$? Discuss.

\item [3.] Compare the expressions for the efficient level of output, $Y^e_t$, and the natural level of output, $Y^n_t$, by taking their ratio. Do they align? If not, explain why. 

\textit{[Hint: Your comparison should be based on two key variables: the markup and the government spending share.]}

\item[4.] Assume now that in the steady state $Y^n=Y^e$. As shown from your answer to point 3, this requires $1-\gamma=1/\mathcal{M}$ or, equivalently, $\gamma=1/\varepsilon$. 

Define two alternative measures of the output gap:
\begin{equation}
x_{t}\equiv \log Y_t-\log Y^e_t=y_{t}-y_{t}^{e}  \label{xe},
\end{equation}
\begin{equation}
\tilde{y_{t}}\equiv \log Y_t-\log Y^e_t=y_{t}-y_{t}^{n}\label{yn}.
\end{equation}

Let a welfare-based loss function of the central bank be of the general form:
\begin{equation}
\mathbb{E}_{0} \sum_{t=0}^{\infty }\beta ^{t}\left \{ \pi
_{t}^{2}+\omega \left( \text{\textit{output gap}}_{t}\right) ^{2}
\right \},  \label{loss}
\end{equation}

subject to a Phillips curve constraint given by:
\begin{equation}
\pi _{t}=\beta \mathbb{E}_{t}\pi _{t+1} + \kappa  \left( \text{\textit{output
gap}}_{t} \right).  \label{nkpc2}
\end{equation}


Which measure of the output gap should appear in the Phillips curve, (\ref{xe}) or (\ref{yn})? Which one should be used in the loss function? Provide an explanation.

%\textit{[Hint: For the rest of the exercise, assume the presence of a wage-bill subsidy, meaning we are analyzing an undistorted steady state.]}

\item[5.] Derive an expression for the gap, if any, between the efficient and natural levels of output, 
$u_{t}$, defined as $u_{t}\equiv \kappa (y_{t}^{e}-y_{t}^{n})$. Determine whether a shock to government spending (to $\gamma _{t}$) can be interpreted as a cost-push shock.

\item[6.] Characterize the effects of a cost-push shock, $u_t$, under the optimal discretionary policy. Is there a stabilization bias? Is the policy time-consistent? 


\end{enumerate}

%%%%%%%%%%%%%%%%%%%%%%%%%%%%%%%%%%%%%%%%%%%
%%%%%%%%%%%%%%%%%%%%%%%%%%%%%%%%%%%%%%%%%%%


%\include{2025/PS/PS 4/ex5_sol}


%%%%%%%%%%%%%%%%%%%%%%%%%%%%%%%%%%%%%%%%%%%
%%%%%%%%%%%%%%%%%%%%%%%%%%%%%%%%%%%%%%%%%%%
%%%%%%%%%%%%%%%%%%%%%%%%%%%%%%%%%%%%%%%%%%%



\end{document}




